\documentclass[a4paper,11pt]{article}
\usepackage{stmaryrd}
\usepackage{graphicx}
\usepackage{amsfonts}
\usepackage{amsmath,amsthm}
\usepackage{amssymb}
\usepackage{mathrsfs}
\usepackage{indentfirst}
\usepackage{titlesec}
\usepackage{caption2}
\usepackage{color}
\usepackage[normalem]{ulem}
\usepackage{algorithm}
\usepackage{algorithmic}
\usepackage{lineno}
\numberwithin{equation}{section}

\usepackage[english]{babel}
\usepackage{cite}
\usepackage{bbm}

%\usepackage[colorlinks, linkcolor=blue]{hyperref}
\usepackage[hypertexnames=false]{hyperref}
%\usepackage[notref,notcite]{showkeys}

\renewcommand{\baselinestretch}{1.2}

\topmargin 0cm \oddsidemargin 0.66cm \evensidemargin 0.66cm
\textwidth 14.66cm \textheight 22.23cm

\parindent 5ex

\providecommand{\bwd}{\leftarrow}
\providecommand{\fwd}{\rightarrow}
\renewcommand{\abstractname}{Abstract}
\renewcommand{\figurename}{Fig.}

\newtheorem{theorem}{Theorem}
\newtheorem{problem}{Problem}
\newtheorem{definition}{Definition}
\newtheorem{condition}{Condition}
\newtheorem{assumption}{Assumption}
\newtheorem{proposition}{Proposition}[section]
\newtheorem{remark}{Remark}
\newtheorem{lemma}{Lemma}[section]
\newtheorem{corollary}{Corollary}[section]
\newtheorem{example}{Example}[section]

%\numberwithin{equation}{section}

%below only needed in Ubuntu 16.04 when using TexStudio, should remove at WinEdt
\usepackage{etoolbox}
\makeatletter
\patchcmd{\ttlh@hang}{\parindent\z@}{\parindent\z@\leavevmode}{}{}
\patchcmd{\ttlh@hang}{\noindent}{}{}{}
\makeatother
%above only needed in Ubuntu 16.04 when using TexStudio, should remove at WinEdt

\usepackage{macros-Wenqing}

\begin{document}

\begin{center}
\textbf{Note on \verb~gpt.py~ computational workflow.}
\end{center}

$\bullet$ RMS norm: normalization of token's embedding vector $\boldx \in  \R^d$. This is different from transformer paper's LayerNorm.

$$\text{RMS Norm}(\boldx)=\dfrac{\boldx}{\text{RMS}(\boldx)+\ve} \ , \ \text{RMS}(\boldx)=\sqrt{\sum\li_{i=1}^d x_i^2} \ , \ \boldx = (x_1,...,x_d) \ , \ \ve>0 \ .$$

\

$\bullet$ Rotary Positional Embedding (RoPE): positional embedding using $\sin$, $\cos$ positional tensors. This is different from transformer paper's Positional Embedding.

Input $x$ with shape $(B, T, H, D)$ where  $B=$batch size, $T=$sequence length, $H=$number of heads in multi-head attention, $D=$dimension of each head (must be even). Let $\sin=(\sin^{t,i})_{0\leq t \leq T-1, 0\leq i\leq D/2}$, $\cos=(\cos^{t,i})_{0\leq t \leq T-1, 0\leq i\leq D/2}$ be tensors each with shape $(1,T,1,D/2)$. Then for each batch $b$, each index $t\in \{0,...,T-1\}$, each head $h$, RoPE does the following:
$$(x^{b,t,h,0},...,x^{b,t,h,D-1}) \ra \left\{\begin{array}{ll}
\cos^{t,i}x^{b,t,h,i} + \sin^{t,i}x^{b,t,h,i+D/2} & \text{ for } 0\leq i \leq D/2-1 \ ,
\\
-\sin^{t,i-D/2}x^{b,t,h,i-D/2} + \cos^{t,i-D/2}x^{b,t,h,i} & \text{ for } D/2\leq i \leq D-1 \ .
\end{array}\right.$$
Position information, originally in $t$, is encoded in such rotation via $\cos, \sin$. Usually we take $\cos^{t,i}=\cos \left(\dfrac{t}{10000^{2i/D}}\right)$ and $\sin^{t,i}=\sin \left(\dfrac{t}{10000^{2i/D}}\right)$.

\

$\bullet$ Causal Self-Attention. The self-attention module with specific causal mask construction.

--  

\end{document}
